\section{Conclusion}


Smart contract auditing is not a single technique---it is a pipeline. Static analysis
catches known patterns. Fuzzing catches implementation edge cases. Symbolic execution
proves invariants. Manual review catches business logic errors. Each phase catches bugs
that the others miss.

Our methodology, applied to 50 contract domains with 832 tests, found 122 issues
including 3 critical vulnerabilities that would have resulted in direct fund loss.
All findings were fixed before deployment. The entire pipeline runs automatically on
every pull request, ensuring that new code does not introduce regressions.

The most important lesson: no tool replaces reading the code. Tools find \emph{known}
bug patterns efficiently. Humans find \emph{unknown} bug patterns---the ones that matter
most.

\begin{thebibliography}{99}
\bibitem{slither} Feist, J. et al. ``Slither: A Static Analysis Framework For Smart Contracts.'' WOOT, 2019.
\bibitem{foundry} Paradigm. ``Foundry: A Blazing Fast, Portable and Modular Toolkit for Ethereum Application Development.'' 2022.
\bibitem{halmos} Park, D. et al. ``Halmos: Symbolic Testing for EVM Smart Contracts.'' a16z Crypto, 2023.
\bibitem{aderyn} Cyfrin. ``Aderyn: Rust-Based Solidity AST Analyzer.'' 2024.
\bibitem{dasp} NCC Group. ``Decentralized Application Security Project (DASP) Top 10.'' 2018.
\bibitem{swc} SmartContractSecurity. ``SWC Registry: Smart Contract Weakness Classification.'' 2020.
\end{thebibliography}

